Notamos que 

\begin{equation*}
   \begin{aligned}
         A_m &= \bar{A}-\bar{B}\bar{m}^T\\
         H & = (\bar{A}-\bar{B}\bar{m}^T)Q\\
            &= \bar{A}Q-\bar{B}\bar{m}^T Q
   \end{aligned}
\end{equation*}

si $G=\bar{m}^T Q$, entonces 

\begin{equation*}
    \begin{aligned}
        H&=\bar{A}Q-\bar{B}G\\
        \bar{m}^T&=GQ^{-1} = GP
    \end{aligned}
\end{equation*}

por lo que , la LMI pasa a ser

\begin{equation*}
    \begin{bmatrix}
        \bar{A}Q-\bar{B}G+(\bar{A}Q-\bar{B}Q)^T+\alpha Q & I \\ I & -\gamma I
    \end{bmatrix} < 0
\end{equation*}

\begin{equation*}
    \begin{bmatrix}
        \bar{A}Q-Q\bar{A}^T+(\bar{B}G-G^T\bar{B})+\alpha Q & I \\ I & -\gamma I
    \end{bmatrix} < 0
\end{equation*}

esta última LMI permite que el readio de la elipsoide sea

\begin{equation*}
    r^2=\frac{\gamma L^2}{\alpha} tr(Q)
\end{equation*}

pues

\begin{equation*}
    \dot{\bar{V}}\leq-\alpha\bar{V}+\gamma L^2
\end{equation*}

\begin{equation*}
    \bar{V}(t)\leq Ve(tr)e^{-\alpha (t-tr)}+ \frac{\gamma L^2}{\alpha}(1-e^{-\alpha (t-tr)})
\end{equation*}

cuanto $t\to \infty$

\begin{equation*}
    \bar{V}(t) \leq \frac{\gamma L^2}{\alpha}
\end{equation*}

con la desigualdad de Rayleigh

\begin{equation*}
    \lambda_{min}(P)||e||^2 \leq \frac{\gamma L^2}{\alpha} = \frac{\gamma L^2 \lambda_{min}(Q)}{\alpha}
\end{equation*}

así, el radio es 

\begin{equation*}
    r \leq \sqrt{\frac{\gamma L^2}{\alpha \lambda_{min}}}= \sqrt{\frac{\gamma L^2 \lambda_{max}(Q)}{\alpha}}
\end{equation*}

los ejes de $||\bar{e}||$ son

\begin{equation*}
    \sum_{i=1}^{n-1} \epsilon_i = \frac{\gamma L^2}{\alpha} tr (Q)
\end{equation*}

de modo que

\begin{equation*}
    ||e||^2 \leq \frac{\gamma L^2}{\alpha} \lambda_{min}(Q) \leq \frac{\gamma L^2}{\alpha} tr(Q)
\end{equation*}

la traza de Q es la suma de los cuadrados de los componentes del error y acotado superiormente por la elipsoide.

\textbf{Ejemplo:} diseñe un controlador SMC-AEM para el sistema 

\begin{equation*}
    \begin{aligned}
        \dot{x}_1&=sin(x_1+x_2)+x_2\\
        \dot{x}_2&=cos(x_1+x_2)+u\\
        Y&=x_1\\
    \end{aligned}
\end{equation*}

